\color{unidarkgreen}{\subsection{ Устройство резервирования отказа выключателя (УРОВ)}}
\color{black}

\begin{enumerate}[label=\arabic{section}.\arabic{subsection}.\arabic*,  align=left, labelsep=4pt, leftmargin=0pt, itemindent=57pt]

\item
Устройство резервирования отказа выключателя (УРОВ) предназначено для выполнения функции
ближнего резервирования защит. В случае отказа выключателя повреждённого присоединения УРОВ
обеспечивает отключение смежных выключателей с целью селективного отделения повреждённого участка
электрической сети.

\begin{figure}[H]
\centering
\includegraphics[width=0.9\textwidth,height=0.9\textheight,keepaspectratio]{\fbpath/661.3311 УРОВ/img/UROV.pdf}
\captionof{figure}{Функциональная схема УРОВ}\label{pict_urov:UROV}
\end{figure}

\item
В устройстве применяется индивидуальный принцип реализации функции УРОВ. Пуск УРОВ происходит
при срабатывании внутренних защит, выявлении неполнофазного режима, либо от внешних защит.

\item
Для выявления факта отказа выключателя предусмотрен контроль наличия тока через выключатель.
Контроль осуществляется измерительным органом УРОВ, порог срабатывания которого задается уставкой
\textbf{«Iср УРОВ»}.

\item
Для обеспечения надежного пуска УРОВ при кратковременном пропадании сигналов пуска до полного
обесточивания отказавшего выключателя предусмотрен подхват от ИО УРОВ, ввод которого осуществляется
с помощью программного ключа \textbf{«Фикс. УРОВ»}, сброс пуска произойдет только при возврате ИО УРОВ.

\item
Возможны два режима работы УРОВ:
\begin{enumerate}[label=\arabic*)]
    \item с автоматической проверкой исправности выключателя;
    \item с дублированным пуском.
\end{enumerate}

\item
В режиме «с автоматической проверкой исправности выключателя» программный ключ
\textbf{«Контр. РПВ»} устанавливается в положение \textit{«Выведено»}, а программный ключ \textbf{«УРОВ на себя»} -- \textit{«Введено»}.
Данный режим выполняет функцию резервирования при обрыве цепи действия на отключение выключателя
от защит. При появлении пуска УРОВ выдается команда на отключение «своего» выключателя, что позволяет
исключить ложное срабатывание УРОВ на смежные выключатели. Выдержка времени действия УРОВ на
«свой» выключатель регулируется уставкой \textbf{«Тср УРОВ В1 на себя»}. Если же отключения выключателя не
происходит, значит неисправна не цепь от защит к ЭМО, а сам ЭМО. Тогда с проверкой срабатывания ИО
УРОВ и с большей выдержкой времени \textbf{«Тср УРОВ»} формируется сигнал срабатывания УРОВ на смежные
выключатели. Контроль наличия тока при действии УРОВ «на себя» может быть выведено программным
ключом \textbf{«Контр I УРОВ на себя»}.

\item
В режиме «с дублированным пуском» программный ключ \textbf{«Контр РПВ»} устанавливается в положение
\textit{«Введено»}, а программный ключ \textbf{«УРОВ на себя»} – \textit{«Выведено»}. В данном режиме пуск УРОВ разрешается
только при срабатывании защит с действием сигнала \textbf{«РПВ от ЭМО»} с действием на отключение --
пропадание входного сигнала ФСУ \textbf{«РПВ от ЭМО»} является признаком срабатывания защит и работе цепи
отключения. С выдержкой времени \textbf{«Тср УРОВ»} формируется сигнал срабатывания УРОВ на смежные
выключатели.

\item
Функциональный блок УРОВ включает в себя логику УРОВ НН, которая выполняется с контролем
положения ШСВ от ЭМВ, пуском от защит (авто)трансформатора, и контролем тока стороны НН
(авто)трансформатора с помощью входных сигналов ФСУ \textit{«Пуск УРОВ НН 1СШ»} и \textit{«Пуск УРОВ НН 2СШ»}.
Применяется в случаях подключения силового (авто)трансформатора по схеме «трансформатор-шины» и
недостаточной чувствительности пусковых органов УРОВ ШСВ при повреждении на стороне НН
(авто)трансформатора. 

\item
Информация о положении выключателя (сигналы «РПВ от ЭМО» и «РПО от ЭМВ») в логике УРОВ
берется непосредственно от соответствующих электромагнитов управления.

\item
Оперативный ввод/вывод УРОВ производится с помощью БУ (блоков управления функциями) \textit{«Вывод УРОВ»}. 

\small
\begin{longtable}{|>{\centering\arraybackslash}m{5.3cm}|>{\centering\arraybackslash}m{3.3cm}|>{\centering\arraybackslash}m{4.2cm}|>{\centering\arraybackslash}m{1.8cm}|>{\centering\arraybackslash}m{1cm}|}
\caption{Перечень уставочных параметров УРОВ\hfill\vspace{-0.5\baselineskip}}\label{desc_urov:tbl1}\\ 
\hline
\rowcolor{gray!30}
Параметр & Условное обозначение на схеме & Значение/ Диапазон & Единица измерения & Шаг \\ 
\hline
\endfirsthead
\caption*{\hspace{3pt}\emph{Продолжение таблицы \ref{desc_urov:tbl1}\hfill\vspace{-0.5\baselineskip}}} \\ % сделано по ГОСТ 2.105 п.6.8.7
\hline
\rowcolor{gray!30}
Параметр & Условное обозначение на схеме & Значение/ Диапазон & Единица измерения & Шаг \\ 
%\hline
\endhead

%===m>УРОВ В1|УРОВ
%===m
\end{longtable}
\normalsize

\end{enumerate}